\begingroup
\fancyhead[L]{\small NACHLASS UND BRIEFWECHSEL}
\clearpage
\SourceLeaf{571}
\pdfbookmark[1]{Nachlass VII. Tafeln für die Störungen der Pallas durch Jupiter}{pallas-nachlass-VII}
\begin{center}
\large NACHLASS UND BRIEFWECHSEL.
\par\bigskip
\Large [VII. TAFELN FÜR DIE STÖRUNGEN\par DER PALLAS DURCH JUPITER.]
\par\medskip
\normalsize [1816 AUGUST---1817 DECEMBER.]
\par\bigskip
[1.]
\par\medskip
\textsc{Gauss} an \textsc{Encke}, Göttingen, 5. August 1816.
\end{center}
\pdfbookmark[2]{[1] Gauss an Encke, 5. August 1816}{pallas-VII-letter1}
\par\noindent\hspace*{\parindent}.\ .\ .\ .\ .\quad Zugleich habe ich noch eine andere Bitte. Die ganze Rechnung
aller Störungen für die Pallas für Eine Opposition ist immer eine Sache von
einem Monat, wegen des Hin- und Herschickens; ich selbst habe (zwar nicht
anhaltend, aber selten hat man dazu auch Lust) diesmal fast 14 Tage an der
einmaligen Rechnung zugebracht. Ich habe die Sache hin und her überlegt,
wie man die Arbeit durch Tafeln abkürzen könnte; und gefunden, dass man
allerdings mit Hülfe einer gewissen Tafel die Arbeit vielleicht in 3 Stunden
abthun könnte. Aber die Berechnung der Tafel wird ungeheure Arbeit
machen. Vor allen Dingen würde nun aber ein Anschlag nöthig sein, wie
\textls[100]{viele} Arbeit, um beurtheilen zu können, ob vielleicht durch 6 oder 8 Rechner
die Sache, wenn auch erst nach Jahren, ausführbar wäre. Zu dem Behuf
schicke ich Ihnen hier die Data zu einer Proberechnung; es liegt mir weniger
an den Resultaten (vorerst wenigstens) als daran, genau zu wissen, wie
viele Zeit Sie etwa dazu gebraucht haben werden. Diese Rechnung möchte
etwa \(\tfrac{1}{100}\) der \textls[100]{ganzen} Arbeit sein.
\par\noindent\hspace*{\parindent}Ich wünsche zu haben die Werthe
\SourceLeaf{572}
\[
\begin{aligned}
\text{1) von }P={}&+359\overset{\prime\prime}{,}834\\
&+54{,}481\cos x &{}+48\overset{\prime\prime}{,}391\sin x\\
&+14{,}295\cos2x &{}-36{,}951\sin2x\\
&+16{,}161\cos3x &{}-8{,}251\sin3x\\
&+6{,}273\cos4x &{}+2{,}923\sin4x\\
&+0{,}162\cos5x &{}+6{,}354\sin5x\\
&-0{,}674\cos6x &{}+3{,}129\sin6x\\
&+1{,}702\cos7x &{}+0{,}899\sin7x\\
&+0{,}922\cos8x &{}-0{,}065\sin8x\\
&+0{,}460\cos9x &{}-0{,}414\sin9x\\
&+0{,}181\cos10x &{}-0{,}326\sin10x\\
&+0{,}063\cos11x &{}-0{,}198\sin11x\\
&+0{,}023\cos12x &{}-0{,}101\sin12x
\end{aligned}
\]
\[
\begin{aligned}
\text{2) von }Q={}&-544\overset{\prime\prime}{,}924\\
&-98{,}084\cos x &{}-6\overset{\prime\prime}{,}178\sin x\\
&+27{,}589\cos2x &{}-10.928\sin2x\\
&+8{,}342\cos3x &{}+4{,}711\sin3x\\
&-0{,}875\cos4x &{}+1{,}294\sin4x\\
&-4{,}480\cos5x &{}-1{,}766\sin5x\\
&-1{,}813\cos6x &{}-1{,}284\sin6x\\
&-0{,}179\cos7x &{}+1{,}643\sin7x\\
&+0{,}434\cos8x &{}+1{,}020\sin8x\\
&+0{,}414\cos9x &{}+0{,}460\sin9x\\
&+0{,}326\cos10x &{}+0{,}181\sin10x\\
&+0{,}198\cos11x &{}+0{,}063\sin11x\\
&+0{,}101\cos12x &{}+0{,}023\sin12x
\end{aligned}
\]
\noindent\hspace*{\parindent}3) von \(\log N\) und \(\psi\), so dass
\[
P=N\cos\psi,\qquad Q=N\sin\psi.
\]
\noindent\(\psi\) wird am Ende auf \textls[100]{Secunden}, \(\log N\) auf 5 Decimalen verlangt. Wünschenswerth
wäre zwar diese Tafel für \(x=0\), \(x=1^0\), \(x=2^0\), \(x=3^0\) etc. bis
\(x=360^0\) zu haben; allein dies wäre eine gar zu furchtbare Arbeit; man
\SourceLeaf{573}
muss sich also mehr beschränken; ich finde dass \textls[100]{allenfalls} man sich begnügen
mag, die Werthe des \(x\) von \(15^0\) zu \(15^0\) wachsen zu lassen, so dass
24 solche Rechnungen zu machen sind. So wünsche ich daher diese Proberechnung
gemacht, wobei übrigens manche Localvortheile zu benutzen sind,
die Sie leicht selbst finden werden und die von
\[
\left.\begin{aligned}
\sin15^0&=\BareRoot{\tfrac{3}{8}}-\BareRoot{\tfrac{1}{8}}\\
\sin45^0&=2\,\mathord{.}\,\BareRoot{\tfrac{1}{8}}
&\qquad\sin60^0&=\BareRoot{\tfrac{3}{4}}\\
\sin75^0&=\BareRoot{\tfrac{3}{8}}+\BareRoot{\tfrac{1}{8}}
\end{aligned}\right\}\ \text{abhangen.}
\]
\noindent\hspace*{\parindent}.\ .\ .\ .\ .
\clearpage
\endgroup
